\makeatletter \def\input@path{{../../}} \makeatother
\documentclass[11pt]{article}

\usepackage[utf8]{inputenc}
\usepackage[T1]{fontenc}
\usepackage[spanish,es-nodecimaldot]{babel}
\usepackage[margin=1in]{geometry}
\usepackage{amsmath,amssymb,mathtools,bm}
\usepackage{enumitem}
\usepackage{fancyhdr}
\usepackage[misc]{ifsym}
\usepackage{microtype}
\usepackage{hyperref}

\setlength\parindent{0pt}
\setlength{\headheight}{14pt}
\pagestyle{fancy}
\lhead{\scshape Pontificia Universidad Cat\'olica de Chile}
\rhead{\scshape IMC UC}
\cfoot{\thepage}
\renewcommand{\headrulewidth}{0.1pt}
\renewcommand{\footrulewidth}{0.1pt}
\renewcommand{\labelitemi}{$\blacklozenge$}

\newcommand{\R}{\mathbb{R}}
\newcommand{\p}{\vspace{8pt}}
\DeclareMathOperator{\diver}{div}
\DeclareMathOperator{\Div}{Div}
\DeclareMathOperator{\tr}{tr}
\DeclareMathOperator{\Cof}{Cof}
\DeclareMathOperator{\sym}{sym}

\begin{document}

\begin{flushleft}
    \small
    \textbf{IMT3130} \hfill
    \textbf{23/05/2026}
\end{flushleft}
\begin{center}
    \LARGE\textbf{Ayudant\'ia 9}\\
    \vspace{3pt}
    \normalsize\textbf{Mec\'anica del continuo: estabilidad, termodin\'amica y linealizaci\'on}\\
    \normalsize Ayudante: Basti\'an Herrera \Letter{\hspace{1pt}: \texttt{biherrera@uc.cl}}
    \rule{\textwidth}{0.05mm}
\end{center}

\section*{Problemas}

\begin{enumerate}[label=\textbf{\arabic*.}, leftmargin=*]

\item \textbf{Stokes--Brinkman II: estabilidad, presi\'on e inf--sup.}

Sea $\Omega\subset\R^d$, $d\in\{2,3\}$, un dominio Lipschitz acotado y conexo. Sean
\begin{equation*}
    V=H^1_0(\Omega)^d,
    \qquad
    Q=L^2_0(\Omega)=\left\{q\in L^2(\Omega):\int_\Omega q\,dx=0\right\}.
\end{equation*}
Se considera el problema mixto de Stokes--Brinkman: hallar $(\bm v,p)\in V\times Q$ tal que
\begin{equation*}
\begin{aligned}
    a(\bm v,\bm w)+b(\bm w,p)&=F(\bm w) &&\forall\bm w\in V,\\
    b(\bm v,q)&=0 &&\forall q\in Q,
\end{aligned}
\end{equation*}
donde
\begin{equation*}
\begin{aligned}
    a(\bm v,\bm w)
    &=2\mu\int_\Omega D(\bm v):D(\bm w)\,dx
      +\alpha\int_\Omega \bm v\cdot\bm w\,dx,\\
    b(\bm w,q)&=-\int_\Omega q\,\diver\bm w\,dx,
\end{aligned}
\end{equation*}
con $\mu>0$, $\alpha\ge0$, y
\begin{equation*}
    D(\bm v)=\frac12(\nabla\bm v+\nabla\bm v^T).
\end{equation*}
Suponga que $F\in V'$ y que se satisface la condici\'on inf--sup continua de la divergencia:
\begin{equation*}
    \inf_{q\in Q\setminus\{0\}}
    \sup_{\bm w\in V\setminus\{\bm0\}}
    \frac{|b(\bm w,q)|}{\|\bm w\|_{H^1(\Omega)}\|q\|_{L^2(\Omega)}}
    \ge \beta>0.
\end{equation*}

\begin{enumerate}[label=(\alph*), leftmargin=*]
    \item Pruebe que $a$ es continua y coerciva en $V$, usando la desigualdad de Korn. Pruebe tambi\'en que $b$ es continua en $V\times Q$.

    \item Considere el caso homog\'eneo $F=0$. Pruebe primero que $\bm v=\bm0$ y luego use la condici\'on inf--sup para probar que $p=0$.

    \item Deduzca la estimaci\'on de estabilidad
    \begin{equation*}
        \|\bm v\|_{H^1(\Omega)}+\|p\|_{L^2(\Omega)}
        \le C\|F\|_{V'}.
    \end{equation*}

    \item Explique por qu\'e el caso $\alpha=0$ sigue estando cubierto, pero el l\'imite puramente Darcy $\mu=0$ no queda cubierto por esta formulaci\'on en $H^1_0(\Omega)^d$.
\end{enumerate}

\item \textbf{Termoelasticidad II: fuente mec\'anica expl\'icita y Lax--Milgram.}

Sea $\Omega_0\subset\R^d$ un dominio material de referencia. Se considera la ecuaci\'on t\'ermica material
\begin{equation*}
    \rho_0c_v\dot\theta-
    \Div_{\bm X}\left(\bm\kappa\nabla_{\bm X}\theta\right)
    =R_0+g_{\mathrm{mec}}
    \qquad\text{en }\Omega_0,
\end{equation*}
con condici\'on adiab\'atica homog\'enea
\begin{equation*}
    \bm\kappa\nabla_{\bm X}\theta\cdot\bm N=0
    \qquad\text{en }\partial\Omega_0.
\end{equation*}
Se asume que $\rho_0>0$, $c_v>0$, y que $\bm\kappa\in L^\infty(\Omega_0)^{d\times d}$ es sim\'etrica y uniformemente el\'iptica:
\begin{equation*}
    \bm\kappa(\bm X)\xi\cdot\xi\ge\kappa_0|\xi|^2
    \qquad\forall \xi\in\R^d,
    \qquad
    \kappa_0>0.
\end{equation*}
La deformaci\'on material es conocida y se denota por
\begin{equation*}
    \bm x=\bm\varphi(\bm X,t),
    \qquad
    F=\nabla_{\bm X}\bm\varphi,
    \qquad
    J=\det F>0.
\end{equation*}
La fuente mec\'anica expl\'icita est\'a dada por
\begin{equation*}
    g_{\mathrm{mec}}
    =3\alpha_TK\theta_\star F^{-T}:\dot F,
\end{equation*}
donde $\alpha_T$, $K$ y $\theta_\star$ se consideran datos conocidos.

\begin{enumerate}[label=(\alph*), leftmargin=*]
    \item Use la identidad de la derivada del determinante para mostrar que
    \begin{equation*}
        F^{-T}:\dot F=\frac{\dot J}{J}.
    \end{equation*}
    Reescriba $g_{\mathrm{mec}}$ en t\'erminos de $J$ y $\dot J$, e interprete brevemente el resultado.

    \item Aplique Euler impl\'icito a la ecuaci\'on t\'ermica y derive la formulaci\'on d\'ebil para calcular $\theta^{n+1}\in H^1(\Omega_0)$, considerando $g_{\mathrm{mec}}^{n+1}$ como dato conocido.

    \item Pruebe existencia y unicidad de $\theta^{n+1}$ usando Lax--Milgram. Indique claramente las hip\'otesis necesarias sobre $R_0^{n+1}$, $g_{\mathrm{mec}}^{n+1}$ y $\theta^n$.

    \item Suponga ahora que, en vez de evaluarse expl\'icitamente, la fuente mec\'anica contiene un t\'ermino de reacci\'on de la forma $m(\bm X)\theta^{n+1}$. Explique qu\'e condici\'on adicional permitir\'ia conservar coercividad.
\end{enumerate}

\item \textbf{Hiperelasticidad finita, principio variacional y linealizaci\'on.}

Sea $\Omega_0\subset\R^d$ un dominio material, con $d\in\{2,3\}$. Consideramos una deformaci\'on
\begin{equation*}
    \bm\varphi(\bm X)=\bm X+\bm u(\bm X),
    \qquad
    F=\nabla_{\bm X}\bm\varphi=I+\nabla_{\bm X}\bm u,
\end{equation*}
y el tensor de Green--Lagrange
\begin{equation*}
    E=\frac12(F^TF-I).
\end{equation*}
Sea un material de Saint Venant--Kirchhoff con energ\'ia almacenada por volumen de referencia
\begin{equation*}
    \Psi(E)=\frac{\lambda}{2}(\tr E)^2+\mu E:E,
\end{equation*}
donde
\begin{equation*}
    \mu>0,
    \qquad
    \lambda+\frac{2\mu}{d}>0.
\end{equation*}
El equilibrio material se escribe
\begin{equation*}
    -\Div_{\bm X}P=\bm B
    \qquad\text{en }\Omega_0,
\end{equation*}
con $\bm u=\bm0$ en $\Gamma_D$, $P\bm N=\overline{\bm T}$ en $\Gamma_N$, y $|\Gamma_D|>0$.

\begin{enumerate}[label=(\alph*), leftmargin=*]
    \item Derive el segundo tensor de Piola--Kirchhoff
    \begin{equation*}
        S=\frac{\partial\Psi}{\partial E},
    \end{equation*}
    y recuerde por qu\'e, para este material hiperl\'astico, el primer tensor de Piola--Kirchhoff viene dado por $P=FS$.

    \item Derive la formulaci\'on d\'ebil lagrangiana no lineal asociada al equilibrio material.

    \item Linealice el modelo alrededor de $F=I$ y obtenga la formulaci\'on d\'ebil de elasticidad lineal.

    \item Pruebe existencia y unicidad del problema linealizado mediante Korn y Lax--Milgram.
\end{enumerate}

\end{enumerate}

\end{document}
